%%%PAGE 109 rot=0 legibility=good summary=简谐运动(回复力判别式_竖直弹簧振子证明F回=-kx_双弹簧_浮体振动_单摆受力与能量分析)
%%%BLOCK id=109-01 ch=08 sec=简谐运动·回复力 kind=knowledge alt=none cont=none y=0.03-0.15
\fbox{简谐运动}\quad $F_{\text{回}}=-kx$　判别式\\[0.8ex]
$a=-\dfrac{kx}{m}=\dfrac{\dd v}{\dd t}$　　$v=\dfrac{\dd x}{\dd t}$　　　振幅与机械能正相关\\[0.8ex]
回复力指向平衡位置（可变方向的\unclear{力}　　变力　不是不变、
%%%END
%%%BLOCK id=109-02 ch=08 sec=简谐运动·弹簧振子 kind=diagram alt=none cont=none y=0.10-0.40
\textbf{①}
%FIG p109_a 0.056 0.112 0.322 0.262
\notefig[0.4]{figures/p109_a.png}
% 图中标注：高 a_max v=0；平 v_max；低 a_max v=0（图中已含）
\[
T=2\pi\sqrt{\dfrac{m}{k}}
\]
%FIG p109_b 0.062 0.287 0.16 0.378
\notefig[0.25]{figures/p109_b.png}
%%%END
%%%BLOCK id=109-03 ch=08 sec=简谐运动·回复力证明 kind=derivation alt=02 cont=none y=0.12-0.40
\textbf{证明}
%FIG p109_c 0.455 0.147 0.646 0.272
\notefig[0.3]{figures/p109_c.png}
\begin{align*}
F_{\text{回}}&=mg+k\Delta x\\
\text{平衡时}\quad mg&=k(x-\Delta x)\\
\Rightarrow F_{\text{回}}&=kx\quad\text{又}\ F_{\text{回}}\text{与}x\text{反向}\\
\text{故}\quad F_{\text{回}}&=-kx
\end{align*}
%FIG p109_d 0.472 0.272 0.646 0.372
\notefig[0.3]{figures/p109_d.png}
\begin{align*}
F_{\text{回}}&=k\Delta x-mg\\
\text{平衡时}\quad mg&=k(\Delta x-x)\\
\therefore F_{\text{回}}&=-kx
\end{align*}
%%%END
%%%BLOCK id=109-04 ch=08 sec=简谐运动·回复力证明 kind=derivation alt=02 cont=none y=0.40-0.62
\textbf{②}　原长一致
%FIG p109_e 0.015 0.405 0.225 0.62
\notefig[0.3]{figures/p109_e.png}
\begin{align*}
F_{\text{回}}&=mg+k_1\Delta x_1+k_2\Delta x_1\\
mg&=k_1(x-\Delta x_1)+k_2(x-\Delta x_1)\\
\therefore F_{\text{回}}&=(k_1+k_2)x
\end{align*}
\underline{又因为$F_{\text{回}}$与$x$方向相反}\\
\struck{又反向}　故$F_{\text{回}}=-(k_1+k_2)x$
%%%END
%%%BLOCK id=109-05 ch=08 sec=简谐运动·对称性 kind=note alt=none cont=none y=0.40-0.46
对称　\fbox{等位移}\ \fbox{等速度}\ \fbox{等加速度}
%%%END
%%%BLOCK id=109-06 ch=08 sec=简谐运动·弹簧振子 kind=derivation alt=06 cont=none y=0.47-0.62
%FIG p109_f 0.645 0.474 0.835 0.585
\notefig[0.25]{figures/p109_f.png}
\begin{align*}
mgh&=\tfrac{1}{2}kh^2\\
mg&=k\Delta x\\
\therefore\Delta x&=\dfrac{h}{2}\\
\therefore h&=2\Delta x
\end{align*}
%%%END
%%%BLOCK id=109-07 ch=08 sec=简谐运动·浮体振动 kind=derivation alt=02 cont=none y=0.63-0.75
\textbf{③}
%FIG p109_g 0.066 0.658 0.185 0.745
\notefig[0.25]{figures/p109_g.png}
\[
F_{\text{回}}=G-F_{\text{浮}}=\rho gS(\Delta x+x)-\rho gS\Delta x=\rho gSx=kx\quad\text{又}\ F\ x\text{反向故}
\] % CHECK: 原稿“G-F浮=ρgS(Δx+x)-ρgSΔx”两项次序与G-F浮的符号不符(应为ρgSΔx-ρgS(Δx+x))，照原样转录
\[
F=-kx
\]
%%%END
%%%BLOCK id=109-08 ch=08 sec=单摆 kind=derivation alt=06 cont=none y=0.75-0.93
\textbf{④}　单摆的力　功　能　位移　加速度、周期分析.
%FIG p109_h 0.024 0.77 0.19 0.925
\notefig[0.3]{figures/p109_h.png}
% 图中标注：θ、a径、Ep max / Ek min、Ek max / Ep min（图中已含）
\[
F_{\text{回}}=mg\sin\theta=mg\dfrac{\widehat{l}}{r}=mg\dfrac{x}{r}=\dfrac{mg}{r}x=kx
\]
$\theta$很小时　$\theta\approx\tan\theta\approx\sin\theta$\\[0.8ex]
振幅$A$与$\bar v_{\text{下}}$，和$T$成正比　　$A=\bar v\dfrac{T}{4}$　　$T=mg\cos\theta+\dfrac{mv^2}{L}$\\[0.8ex]
机械能守恒　　$a_{\text{切}}=g\sin\theta$　　$a_{\text{径}}=\dfrac{V_{\unclear{\mathrm{max}}}^2}{L}=\dfrac{v^2}{L}$
%%%END
%%%PAGEEND 109
